\section*{Complementary spectral gap}
Work in the real Hilbert space
$L^2((0,\pi);\mathbb R^{m+1})$ restricted by the integrated conservation law,
with domain $H_N^2$ subject to the same restriction.  The diagonal diffusion
operator with positive entries is sectorial and has compact resolvent; the
reaction derivative is a bounded multiplication perturbation.

At $\mu=0$, the noncenter spectrum consists of:
\begin{enumerate}
\item the spectrum of $A_m|_{c^\perp}$;
\item the nonzero spectrum of $A_m-D_m$;
\item the spectra of $A_m-k^2D_m$, $k\ge2$.
\end{enumerate}
Every listed eigenvalue has negative real part.  For large $k$, diagonal
diffusion domination sends the spectral bound to $-\infty$; hence only
finitely many modes can lie in any fixed right half-plane.  The discrete
complementary spectrum therefore has a positive gap from the imaginary axis
for each fixed $m$.  The homogeneous conservation zero is absent because the
phase space fixes the integrated mass and $\ker A_m\cap c^\perp=\{0\}$.
